% Image credits for the photographs and AI illustrations of Book 2
% (University Chemistry, Year 1). Machine-readable license records live in
% images/book2/CREDITS.md; keep the two files in sync. The Book 2 writer
% owns this file: add one \item per photograph (in an itemize, after the
% first paragraph) as photographs land.
\chapter*{Créditos das imagens}

Os desenhos deste livro são originais. As fotografias, quando usadas, são
utilizadas sob as respectivas licenças livres, com agradecimento aos seus
autores; os links completos para as fontes e as URLs das licenças de cada
fotografia estão listados em \texttt{images/book2/CREDITS.md} no
repositório do livro.
Os pictogramas de perigo são os do Sistema Globalmente Harmonizado de
Classificação e Rotulagem de Produtos Químicos, tal como desenhados para o
pacote \texttt{ghsystem} do \TeX\ Live.

\bigskip
Ao lado das fotografias e dos desenhos originais, algumas figuras são
ilustrações geradas por IA, produzidas com a geração de imagens da OpenAI a
partir de instruções escritas pelos editores do livro, e revisadas quanto à
exatidão química antes de serem incluídas. A instrução completa de cada
imagem gerada está registrada em \texttt{images/book2/ai/PROMPTS.md} no
repositório.

\bigskip
\noindent Fotografias, por capítulo:
\begin{itemize}\small
  \item Cap.~1: Niels Bohr, por volta de 1922, AB Lagrelius \& Westphal,
        domínio público (Wikimedia Commons).
  \item Cap.~2: Linus Pauling, 1962, Fundação Nobel, domínio público
        (Wikimedia Commons).
  \item Cap.~5: amostra de cobre nativo, de Flipin, OG, CC0 (Wikimedia
        Commons).
  \item Cap.~6: cristais de halita, fluorita e um diamante octaédrico, de
        James St. John, CC BY 2.0; grafite, de Zbynek Burival,
        CC BY-SA 4.0; cristais de neve de Wilson A. Bentley (1902), domínio
        público (todas da Wikimedia Commons).
  \item Cap.~8: Svante Arrhenius, 1909, Meisenbach Riffarth \& Co.,
        domínio público (Wikimedia Commons).
  \item Cap.~11: precipitados de cloreto, brometo e iodeto de prata, de
        Rrausch1974, CC BY-SA 3.0 (Wikimedia Commons).
  \item Cap.~12: suspensão de hidróxido de cobre(II), de Alvy16, CC BY 4.0;
        solução de tetra-aminocobre(II), de Chemicalinterest, domínio público (ambas
        da Wikimedia Commons).
  \item Cap.~13: Walther Nernst, 1889, Smithsonian Libraries, domínio
        público (Wikimedia Commons).
  \item Cap.~16: Jacobus Henricus van 't Hoff, fotografia publicada em
        1911, domínio público (Wikimedia Commons).
  \item Cap.~21: Victor Grignard, domínio público (Wikimedia Commons).
  \item Cap.~25: Vladimir Markovnikov, retrato de um obituário publicado
        em 1905, domínio público (Wikimedia Commons).
  \item Cap.~26: lagoas de lítio do Salar de Atacama, NASA Earth
        Observatory, domínio público; sódio em óleo mineral, de Justin
        Urgitis, CC BY-SA 2.5; cores de chama, de Hegelrast, CC BY-SA 4.0
        (todas da Wikimedia Commons).
  \item Cap.~27: bauxita pisolítica, de James St. John, CC BY 2.0; bórax,
        de Aram Dulyan, domínio público; quartzo, de Stephanie Clifford,
        CC BY 2.0; fósforo branco, Hi-Res Images of Chemical Elements,
        CC BY 3.0; Fritz Haber, Fundação Nobel, domínio público (todas da
        Wikimedia Commons).
  \item Cap.~28: extração de enxofre na cratera do Kawah Ijen, de
        S\'emhur, CC BY-SA 4.0; cloro em uma ampola, de W. Oelen,
        CC BY-SA 3.0; bromo em uma ampola, de Jurii, CC BY 3.0; cristais de
        iodo, de Dnn87, CC BY 3.0 (todas da Wikimedia Commons).
  \item Cap.~29: banco aquecido (banco de Kofler), de HBR, CC BY-SA 3.0;
        refratômetro de Abbe, de Foreade, CC BY-SA 4.0 (ambas da Wikimedia
        Commons).
\end{itemize}
\clearpage
